\section{Companion lab: compute a SHA-256 fingerprint}
\label{revision-lab:sha256}

A ChaCha20 block generates bytes to combine with a message. SHA-256 instead
reduces a message to a fixed 32-byte \emph{digest}, also called a fingerprint.
The calculation must depend on every input byte and on where the message ends.
A useful implementation therefore needs more than a compression loop: it needs
a representation of unfinished input and an explicit finalization operation.
This companion revises original project \texttt{025-sha256}.

Use the lesson package \emph{Carry a message through SHA-256} alongside this
companion. Its short worked programs expose particular boundaries; you complete
the six entries in one continuing \texttt{independent/sha256.asm} source. The book checkout contains the same starter at
\path{revision/labs/files/sha256.asm}.
It supplies the caller, data arenas and constants; each unfinished entry stops
with \texttt{break}. Choose your own private labels, temporaries, helpers and
frames under the public contracts. Functions, file ownership and word
transformations are prerequisites.

Answer the inline predictions, then set a breakpoint on the requested instruction
and Debug that source. The selected paused frame and blue execution row identify
where the program stopped. After accepting an observation, the lesson may already
point to the next instruction in yellow. Move the breakpoint to that next target
and use Resume in the same process. Finish the process after the final stop,
before launching another worked program.

\begin{figure}[!ht]
\centering
\begin{tikzpicture}[alt={Three crops from one actual light-theme IDE photograph. The blue execution row is the store at schedule.asm line 27, confirmed by the selected paused frame. Register t0 contains hexadecimal 0000000061626380. The yellow lesson pointer has already moved to words_ready at line 34, the next requested stop. Label A identifies the current paused instruction; label B identifies the next lesson target.}]
\node[anchor=north west,inner sep=0,font=\small\bfseries] (title) at (0,0) {Distinguish the pause from the next lesson target};
\node[anchor=north west,inner sep=0] (code) at ([yshift=-2mm]title.south west) {\MARSRevisionCrop{revision/captures/sha/sha-native-pause.png}{0.96\linewidth}{0.42968750}{0.53666667}{0.17968750}{0.23888889}};
\begin{scope}[shift={(code.south west)},x={($(code.south east)-(code.south west)$)},y={($(code.north west)-(code.south west)$)}]
\draw[black,line width=.7pt] (.11,.7723) rectangle (.995,.8960);
\draw[black,line width=.7pt] (.11,.01485) rectangle (.995,.13366);
\node[draw=black,fill=white,inner sep=1pt,font=\small\bfseries] at (.035,.84) {A};
\node[draw=black,fill=white,inner sep=1pt,font=\small\bfseries] at (.035,.075) {B};
\end{scope}
\node[anchor=north west,inner sep=0] (frame) at ([yshift=-2mm]code.south west) {\MARSRevisionCrop{revision/captures/sha/sha-native-pause.png}{0.4416\linewidth}{0.03750000}{0.21555556}{0.78281250}{0.74888889}};
\node[anchor=north west,inner sep=0] (value) at ([yshift=-2mm]frame.south west) {\MARSRevisionCrop{revision/captures/sha/sha-native-pause.png}{0.6528\linewidth}{0.24765625}{0.07666667}{0.48671875}{0.89777778}};
\end{tikzpicture}
\caption{The selected frame confirms the current pause at line 27 (A), before
\texttt{sw} stores numeric word \texttt{61626380}. The lesson has accepted that
observation and points to line 34 (B). Move the breakpoint to the next requested
instruction and use Resume. The frame and register crops belong to the same
photograph; the program has not advanced to the yellow target.}
\label{fig:revision-sha-native-pause}
\end{figure}


The implementation checks follow the dependencies: schedule, compression,
initialization, update, finalization and file ownership. The update-only checks
inspect a returned context before finalization is available. Construction edits
can invalidate earlier results, so the last assessment rechecks all original
cases on the completed source. Keep a reasoning argument beside those results;
a finite test grid does not explain why another input should work.

The algorithm is specified in
\href{https://doi.org/10.6028/NIST.FIPS.180-4}{FIPS 180-4}.
\href{https://www.rfc-editor.org/rfc/rfc6234.html}{RFC 6234} provides additional
implementation and test material. Use those references for the fixed definitions,
not as a substitute for tracing your own state transitions. The local assignment
specifies our context layout, bounded calls and simulator file interface.

\subsection*{Distinguish input bytes, numeric words and printed digits}

Start with the three ASCII bytes for \texttt{abc}: hexadecimal
\texttt{61 62 63}. No terminating zero or newline is part of this message.
Its digest has 32 bytes. Printing each byte as two hexadecimal digits produces
64 characters; those characters are a representation of the digest, not another
32-byte message to hash.

Internally, SHA-256 processes 64-byte blocks. Its first sixteen schedule words
interpret successive groups of four block bytes in big-endian order.
The bytes \texttt{61 62 63 80} therefore represent numeric word
\texttt{61626380}. Loading that group with a native little-endian word load
would instead produce \texttt{80636261}. Decode the bytes deliberately, even
though the block happens to lie at an aligned address in one example.
The interface also accepts unaligned block starts.

After decoding, the schedule is an array of \emph{numeric} words. A normal word
store in this simulator represents \texttt{61626380} in memory as
\texttt{80 63 62 61}. That storage order is correct for the schedule array.
The final digest has a different interface: it must be written as big-endian
bytes. Decide what an object represents before deciding whether its byte order
is wrong. A byte view and a word view can both be correct.

Implement \texttt{sha256\_schedule(a0 block,a1 schedule)} first. It reads
exactly 64 block bytes and replaces all 64 words of a word-aligned, 256-byte
schedule. Previous workspace contents must not affect the result. Source and
output are disjoint. Run \texttt{checks/kernel.kts} while developing; failures
for entries you have not implemented are expected. Use the individual case names
to distinguish a completed schedule from incomplete compression or streaming.

\subsection*{Expand a schedule with four explicit predecessors}

The first sixteen words come directly from the block. For indices 16 through
63, compute a wrapping word sum:
\[
 W_t=W_{t-16}+\sigma_0(W_{t-15})+W_{t-7}+\sigma_1(W_{t-2})
 \pmod {2^{32}}.
\]
The two small-sigma functions combine rotations and a logical shift:
\[
\begin{aligned}
\sigma_0(x)&=\operatorname{rotr}(x,7)\mathbin{\oplus}
 \operatorname{rotr}(x,18)\mathbin{\oplus}(x\gg3),\\
\sigma_1(x)&=\operatorname{rotr}(x,17)\mathbin{\oplus}
 \operatorname{rotr}(x,19)\mathbin{\oplus}(x\gg10).
\end{aligned}
\]
Here the right shifts insert zeros. An arithmetic shift would copy the sign bit
and change the function for words whose high bit is one. All rotations and sums
operate on 32-bit words, just as in the quarter-round chapter.

If a pointer names the destination for $W_t$, then its four predecessors are
at byte offsets $-64$, $-60$, $-28$ and $-8$. Derive each offset by multiplying
the index difference by four. This gives a way to debug an incorrect recurrence
without guessing at constants.

\MARSCodeLabel{One possible allocation for the first small-sigma term}
\begin{minted}{mars-idea}
    lw $t2, -60($t0)
    rotr $t3, $t2, 7
    rotr $t4, $t2, 18
    xor $t3, $t3, $t4
    srl $t4, $t2, 3
    xor $t3, $t3, $t4
\end{minted}

This fragment assumes \texttt{t0} points to the next schedule word, not to the
start of the schedule. It retains the unmodified input in \texttt{t2} until
both rotations and the shift are complete. A different allocation is equally
valid if it preserves that dependency.

For the padded \texttt{abc} block developed below, $W_0$ is
\texttt{61626380}, $W_1$ through $W_{14}$ are zero, and $W_{15}$ is
\texttt{00000018}. The next four words are \texttt{61626380},
\texttt{000f0000}, \texttt{7da86405} and \texttt{600003c6}.
Predict $W_{16}$ using its four predecessors before running. Then inspect all
four expanded words. If $W_{16}$ is correct and $W_{17}$ is wrong, examine which
predecessors changed between those two iterations.

\begin{figure}[!ht]
\centering
\begin{tikzpicture}[alt={Actual native Memory row for schedule words sixteen through nineteen. Bytes are 80 63 62 61, 00 00 0f 00, 05 64 a8 7d and c6 03 00 60. An outline identifies the first numeric word.}]
\node[anchor=north west,inner sep=0,font=\small\bfseries] (title) at (0,0) {Expanded schedule words use native storage};
\node[anchor=north west,inner sep=0] (a) at ([yshift=-2mm]title.south west) {\MARSRevisionCrop{revision/captures/sha/sha-schedule.png}{0.99\linewidth}{0.03125}{0.33888889}{0.58203125}{0.60888889}};
\begin{scope}[shift={(a.south west)},x={($(a.south east)-(a.south west)$)},y={($(a.north west)-(a.south west)$)}]
\draw[black,line width=.7pt] (.20,.01) rectangle (.39,.467);
\end{scope}
\end{tikzpicture}
\caption{The pictured row contains numeric words W16 through W19. Its first four bytes, 80 63 62 61, store word 61626380 in native little-endian order. An outline identifies that first word.}
\label{fig:revision-sha-native-schedule}
\end{figure}


\subsection*{Keep a round's old state until its new values are ready}

\texttt{sha256\_compress(a0 state,a1 block,a2 workspace)} receives eight
aligned numeric state words, one 64-byte block and a 256-byte workspace.
It expands that block's schedule, performs 64 rounds and adds the working
results to the original state words. The source block is preserved. This is one
compression operation; it does not choose the initial state or add padding.

Name the eight working words $a$ through $h$. The large-sigma functions contain
three rotations each:
\[
\begin{aligned}
\Sigma_0(x)&=\operatorname{rotr}(x,2)\mathbin{\oplus}
 \operatorname{rotr}(x,13)\mathbin{\oplus}\operatorname{rotr}(x,22),\\
\Sigma_1(x)&=\operatorname{rotr}(x,6)\mathbin{\oplus}
 \operatorname{rotr}(x,11)\mathbin{\oplus}\operatorname{rotr}(x,25).
\end{aligned}
\]
They are different functions from the small sigmas in the schedule.
Two further functions operate independently at each bit position.
\emph{Choice}, $\operatorname{Ch}(e,f,g)$, selects the bit from $f$ when the
corresponding bit of $e$ is one, and from $g$ otherwise. \emph{Majority},
$\operatorname{Maj}(a,b,c)$, is one when at least two of those three input bits
are one. They can be expressed as
$(e\mathbin{\&}f)\mathbin{\oplus}(\mathord{\sim}e\mathbin{\&}g)$ and
$(a\mathbin{\&}b)\mathbin{\oplus}(a\mathbin{\&}c)
\mathbin{\oplus}(b\mathbin{\&}c)$, respectively, taking complements within
the word.

At round $t$, form two temporary word sums from the \emph{old} working state:
\[
\begin{aligned}
T_1&=h+\Sigma_1(e)+\operatorname{Ch}(e,f,g)+K_t+W_t,\\
T_2&=\Sigma_0(a)+\operatorname{Maj}(a,b,c).
\end{aligned}
\]
The supplied table gives the fixed $K_t$ values; do not derive them from message
bytes. The new working sequence is
$(T_1+T_2,a,b,c,d+T_1,e,f,g)$, with all sums modulo $2^{32}$.
Both new $a$ and new $e$ need $T_1$, and new $e$ needs old $d$.

\Needspace{15\baselineskip}
For example, if \texttt{s0} through \texttt{s7} hold $a$ through $h$, and
\texttt{t0/t1} hold $T_1/T_2$, the following order avoids destroying a value
before its last use:

\MARSCodeLabel{Publish one round's eight new working words}
\begin{minted}{mars-idea}
    move $s7, $s6
    move $s6, $s5
    move $s5, $s4
    addu $s4, $s3, $t0
    move $s3, $s2
    move $s2, $s1
    move $s1, $s0
    addu $s0, $t0, $t1
\end{minted}

Trace why reversing these moves fails. Moving $a$ into $b$ before preserving old
$b$ for $c$, for example, duplicates $a$ instead of shifting the old state.
Using saved registers for working words also means the function must save and
restore the caller's complete 64-bit values. The fact that the calculation uses
only words does not narrow the calling convention.

Starting with the supplied SHA-256 initial values and the padded \texttt{abc}
block, the first round gives $T_1=$\texttt{54da50e8} and
$T_2=$\texttt{08909ae5}. New $a$ is \texttt{5d6aebcd}; new $e$ is
\texttt{fa2a4622}. Use this early boundary to separate a wrong Boolean function
or constant from an error later in the loop. After round 64, add each working
word to its corresponding \emph{original} state word. Storing only the working
sequence omits feed-forward.

The next block of the same message starts from this updated chaining state.
Its rounds therefore depend on the preceding block. The independent columns
that made the ChaCha20 MSA arrangement useful are absent here. Several messages
could be processed independently, but that would require keeping a separate
context for each message.

\subsection*{Represent unfinished input in the context}

A caller should be able to supply \texttt{abc} at once or as updates containing
\texttt{a} and \texttt{bc}. The digest must be the same. Padding each update
would insert extra bytes between those chunks, so an update cannot finalize the
message merely because its current input ends.

\begin{table}[!ht]
  \centering
  \begin{tabular}{r r l}
    \toprule
    Byte offset & Bytes & Meaning                                    \\
    \midrule
    0           & 32    & Eight numeric chaining words               \\
    32          & 8     & Original message byte total                \\
    40          & 4     & Used bytes in the unfinished block         \\
    44          & 4     & Reserved                                   \\
    48          & 64    & Block buffer; only the used prefix is live \\
    112         & 256   & Schedule workspace                         \\
    \bottomrule
  \end{tabular}
  \caption{The 368-byte context. Numeric fields use the simulator's native little-endian storage. The source block and published digest have their own big-endian interpretation. Unused buffer and workspace bytes need not be zero.}
  \label{tab:revision-sha-context}
\end{table}


\texttt{sha256\_init(a0 context)} installs the eight standard initial words
and sets total and used to zero independently of old contents. It need not clear
unused storage. \texttt{sha256\_update(a0 context,a1 source,a2 count)} accepts
counts zero through 4096. Total original message bytes remain below $2^{61}$.
The context, source and other supplied objects are disjoint.

\Needspace{7\baselineskip}
On return from update, the chaining words describe all complete message blocks,
used equals total modulo 64, and the live buffer prefix contains the remaining
original bytes. These three facts are the invariant. Once used reaches 64,
compress the buffer and reset used to zero. Buffer bytes beyond the new used
prefix are stale storage, not additional input.

For a 65-byte message supplied as 63 bytes and then two, the first call leaves
used 63. The next byte completes a block; compress it. The final byte becomes
buffer byte zero, leaving total 65 and used one. Before coding, trace the source
pointer, remaining count and used value at all three points. Do not let a nested
compression call destroy the outer source walker.

An empty update reads no source bytes, even if its source pointer is deliberately
invalid. It can still inspect its valid context. Calling update on two contexts
in alternating order must preserve each one's independent state and tail.
Shared read-only constants are safe; one shared mutable buffer is not.

In \texttt{worked/fill-block.asm}, the two-byte request increases the original
byte total to 65 once, then only one byte is copied before the \texttt{block\_ready}
pause. Used is temporarily 64 and one input byte remains. The state words have
not absorbed the full block. This limited program stops at an internal boundary;
your update must perform that compression and finish the remaining fragment
before returning a valid context.

With mode zero, the supplied caller uses three original bytes. After update,
inspect total three, used three and the live prefix \texttt{61 62 63} before
finalization.

\subsection*{Prove whether padding needs one block or two}

\texttt{sha256\_final(a0 context,a1 output32)} consumes a prepared context
and publishes exactly 32 bytes. Finalization appends a single one bit, followed
by enough zero bits and the original message's 64-bit bit length. Since our
messages contain whole bytes, the marker byte is \texttt{80}. It is not a
terminator to search for in the original message; an original \texttt{80} byte
is ordinary data.

The final eight bytes of the last block hold the length, most significant byte
first. For \texttt{abc}, the marker is at buffer offset three, zeros extend
through offset 55, and offsets 56--63 contain
\texttt{00 00 00 00 00 00 00 18}. The final byte is hexadecimal 18 because
three original bytes contain 24 bits.

\begin{table}[!ht]
\centering
\begin{tabular}{r r r r}
\toprule
Original bytes & Blocks handled by update & Buffered bytes & Final blocks\\
\midrule
0 & 0 & 0 & 1\\
3 & 0 & 3 & 1\\
55 & 0 & 55 & 1\\
56 & 0 & 56 & 2\\
63 & 0 & 63 & 2\\
64 & 1 & 0 & 1\\
65 & 1 & 1 & 1\\
\bottomrule
\end{tabular}
\caption{Finalization always processes padding, including after an exact full message block. The last column counts additional compression calls during finalization, not the total calls for the message.}
\label{tab:revision-sha-padding}
\end{table}


\Needspace{9\baselineskip}
With 55 buffered bytes, the marker fits at offset 55 and the length starts at
56 in the same block. With 56 buffered bytes, the marker occupies offset 56,
where the length would need to start. Zero the rest of that block and compress
it first. Then create a second block with zeros through offset 55 and the
length in the last eight bytes. A full original 64-byte block has already been
compressed by update, but finalization still needs a padding block.

\begin{figure}[!ht]
\centering
\begin{tikzpicture}[alt={Actual native Memory view of the last thirty-two bytes of the second padding block for a fifty-six-byte message. All pictured bytes are zero except the final two, 01 c0. An outline identifies the eight-byte big-endian original bit length.}]
\node[anchor=north west,inner sep=0,font=\small\bfseries] (title) at (0,0) {The final block retains the original bit length};
\node[anchor=north west,inner sep=0] (a) at ([yshift=-2mm]title.south west) {\MARSRevisionCrop{revision/captures/sha/sha-padding.png}{0.99\linewidth}{0.03125}{0.31555556}{0.58203125}{0.60888889}};
\begin{scope}[shift={(a.south west)},x={($(a.south east)-(a.south west)$)},y={($(a.north west)-(a.south west)$)}]
\draw[black,line width=.7pt] (.58,.005) rectangle (.965,.303);
\end{scope}
\end{tikzpicture}
\caption{For 56 original bytes, the last two rows of the second padding block are shown before its compression. The outlined eight-byte field ends 01 c0, encoding 448 bits. Earlier message bytes and the marker are outside this crop.}
\label{fig:revision-sha-native-padding}
\end{figure}


\Needspace{8\baselineskip}
Compute the bit length from the original byte total using a doubleword left
shift by three. Serialize all eight bytes with full-width shifts. A word shift
cannot extract the upper half correctly. The public
\texttt{worked/padding-high-length.asm} experiment supplies a synthetic original
byte total of $2^{32}+63$ and a 63-byte unfinished tail. It constructs two final
blocks whose length field is \texttt{00 00 00 08 00 00 01 f8}. These values let
you inspect the high length bits without reading a file of that size. The
experiment constructs padding bytes only; the actual compression belongs to
your finalization function.

After the final compression, serialize each of the eight numeric state words
most significant byte first. For \texttt{abc}, the first four digest bytes are
\texttt{ba 78 16 bf}. Storing numeric word \texttt{ba7816bf} with a native
word store would reverse those four output bytes. Check the complete digest,
not only this prefix, against an independent digest implementation.

The \texttt{worked/serialize.asm} experiment is given the final numeric state
for \texttt{abc}. It demonstrates this output conversion without calculating
that state. Your compression and finalization must obtain the correct words
before applying the same byte-order rule.

\begin{figure}[!ht]
  \centering
  \begin{tikzpicture}[alt={Actual native Memory view of all thirty-two SHA-256 digest bytes for abc. The first row begins ba 78 16 bf and the second begins b0 03 61 a3. An outline identifies the first four serialized bytes.}]
    \node[anchor=north west,inner sep=0,font=\small\bfseries] (title) at (0,0) {The digest is a serialized byte sequence};
    \node[anchor=north west,inner sep=0] (a) at ([yshift=-2mm]title.south west) {\MARSRevisionCrop{revision/captures/sha/sha-digest.png}{0.99\linewidth}{0.03125}{0.31555556}{0.58203125}{0.60888889}};
    \begin{scope}[shift={(a.south west)},x={($(a.south east)-(a.south west)$)},y={($(a.north west)-(a.south west)$)}]
      \draw[black,line width=.7pt] (.20,.3088) rectangle (.39,.6176);
    \end{scope}
  \end{tikzpicture}
  \caption{All 32 digest bytes for the three-byte message abc are shown. The outlined prefix is ba 78 16 bf. Interpret the Bytes column in order; a native little-endian word view groups these digest bytes differently.}
  \label{fig:revision-sha-native-digest}
\end{figure}


\Needspace{6\baselineskip}
Finalization consumes the context. Its contents afterward are unspecified,
and calling final twice is outside the interface. Initialize it again before
hashing another message. Do not add padding bytes to the saved original-byte
count while calculating their own length field.

\subsection*{Finalize a file only after successful EOF}

The file function receives a terminated path, scratch pointer, capacity from one
through 256, and a 32-byte digest destination. It owns a fresh private context.
It opens the path read-only, initializes that context and repeatedly reads a
positive request no larger than capacity.

A positive result tells how many bytes are now available in scratch. Supply
exactly that count to update. A short positive result does not mean EOF; the next
read may produce more data. A negative result is a read error. Zero is successful
EOF, at which point finalization can publish the digest.

Suppose capacity is seven but a controlled reader returns three bytes, then two,
then minus one. The context has consumed five successful bytes, but there is no
complete file digest to publish. Return status two, close the acquired descriptor
and preserve all 32 destination bytes. Hashing the requested seven bytes would
incorporate stale scratch contents; treating minus one as EOF would publish a
digest of a prefix.

The public \texttt{worked/read-prefixes.asm} program separates file ownership
from hash arithmetic. It counts successful input bytes and publishes one count
only at EOF. Its successful lesson fixture returns two bytes, then one, then
EOF; the error fixture returns two bytes, then minus one, preserving the original
output word. The SHA wrapper adds a fresh context, updates and finalization
under the same success/failure ownership rule, applied to all 32 digest bytes.

\begin{figure}[!ht]
  \centering
  \begin{tikzpicture}[alt={Actual native Memory view after a deterministic file read error. Both rows of the thirty-two-byte digest destination remain a5. An outline encloses the unchanged destination. Status and descriptor closure are checked separately, not visible in the crop.}]
    \node[anchor=north west,inner sep=0,font=\small\bfseries] (title) at (0,0) {A failed read leaves the digest unpublished};
    \node[anchor=north west,inner sep=0] (a) at ([yshift=-2mm]title.south west) {\MARSRevisionCrop{revision/captures/sha/sha-read-error.png}{0.99\linewidth}{0.03125}{0.31555556}{0.58203125}{0.60888889}};
    \begin{scope}[shift={(a.south west)},x={($(a.south east)-(a.south west)$)},y={($(a.north west)-(a.south west)$)}]
      \draw[black,line width=.7pt] (.20,.005) rectangle (.965,.6176);
    \end{scope}
  \end{tikzpicture}
  \caption{After successful reads of three and two bytes followed by a read error, the full 32-byte destination remains a5. The outline encloses that unchanged destination. Separate checks verify status two and closing descriptor 19; those values are not pictured.}
  \label{fig:revision-sha-native-read-error}
\end{figure}


An open failure returns status one and acquires no descriptor. On success and
read failure, close the descriptor actually returned by open, which need not be
three. Keep that descriptor separate from argument registers reused by update
and finalization. Also keep the result status across the close syscall so that
its return value does not replace the function's intended status.

\Needspace{12\baselineskip}
With caller defaults retained (\texttt{mode} two, \texttt{length} three and
\texttt{buffer\_capacity} seven), prepare an ordinary host-file observation from
your completed source. From the book checkout, choose a fresh output path:

\Needspace{7\baselineskip}
\begin{Verbatim}[fontsize=\small]
python3 packages/025-sha256/prepare.py s45-1 \
  --source path/to/your/main.asm \
  --output /private/tmp/sha-file-observation.asm \
  --fixture-output /private/tmp/sha-input.bin
\end{Verbatim}

The tool refuses existing outputs. The resulting file contains deterministic
binary bytes and the source names that exact host path. An ordinary file run
checks real file integration; it does not force a short read or injected error.
The public file checks supply those separate controls.

\subsection*{Build evidence at the boundary that can fail}

Complete the independent source in place and retain the worked programs for
separate observations. For optional standalone data experiments, the preparation
tool can apply \texttt{s03-1} for the padded schedule input, \texttt{s11-1} for a
compression input with the standard initial state, \texttt{s28-2} for crossing
a buffered block boundary, and \texttt{s37-3} for two-block finalization.
These recipes install data and controls while retaining your instruction text.
Do not reuse a context after finalization just because its old buffer still
looks plausible.

Implement and test schedule, compression, initialization, update, finalization
and file ownership in that dependency order. The public completed suites are
\texttt{kernel.kts}, \texttt{file.kts}, \texttt{program.kts} and
\texttt{reuse.kts} under the package's \texttt{checks} directory. Run each with
\texttt{java -jar /path/to/cli.jar grade CHECK SOURCE}.
The caller stops at \texttt{main\_done} before storing its returned status and
at \texttt{results\_ready} after the store. These are different observation
boundaries; the later exit syscall changes \texttt{v0}.

Write a test plan covering 0, 3, 55, 56, 63, 64 and 65 bytes, then binary input
containing zeros and bytes above 127. For each length, compare one update with
several splits, including empty updates. Interleave two contexts and initialize
a used context again. Move the selected slices within their supplied arenas.
A correct digest at a fixed address does not establish correct pointer use or
memory ownership.

For error diagnosis, check the first wrong schedule word, first wrong round,
first wrong context invariant or first wrong serialized byte. A complete
fingerprint is valuable as an end-to-end check, but it cannot identify which
of those boundaries first diverged.

Finally, an unkeyed digest does not identify a sender. A party that can replace
both a message and its advertised digest can recompute them together. Keyed
message authentication requires a defined construction, such as \href{https://www.rfc-editor.org/rfc/rfc2104.html}{HMAC}; appending
a secret to this exercise is not its specification. The work here is a precise
bounded SHA-256 implementation with reproducible observations. Explain its
state and ownership promises alongside the passing results.
